% f3_orient.tex -- Figure 3: the seven other orientations of the square, and what they do to the amplitudes (natural width about 17.2 cm).
%   (a) the eight orientations of one CIFAR-10-like image (bird.png of the slides): as it is (navy frame) and the seven G_k.
%   (b) a schematic of the per-block mean amplitudes (gen/mk_data.py, not data): ID (navy, band = +-1 SD), the seven re-oriented means
%       (grey), an OOD mean (red).  The measured shares are in the text (main resume, Table 7).
\begin{tikzpicture}[>={Stealth[length=1.8mm, width=1.4mm]}, line cap=round, line join=round, font=\footnotesize,
  lab/.style={inner sep=1pt, align=center, font=\footnotesize}, sm/.style={inner sep=1pt, align=center, font=\scriptsize},
  ttl/.style={anchor=west, inner sep=0pt, font=\small}]
% written by gen/mk_data.py: schematic per-block mean amplitudes (drawing coordinates, cm), p = 23
\def\profid{(8.350,-1.893) -- (8.716,-1.897) -- (9.082,-1.818) -- (9.448,-1.685) -- (9.814,-1.676) -- (10.180,-1.858) -- (10.545,-2.081) -- (10.911,-2.175) -- (11.277,-2.176) -- (11.643,-2.238) -- (12.009,-2.385) -- (12.375,-2.452) -- (12.741,-2.322) -- (13.107,-2.096) -- (13.473,-1.965) -- (13.839,-1.952) -- (14.205,-1.908) -- (14.570,-1.760) -- (14.936,-1.642) -- (15.302,-1.712) -- (15.668,-1.924) -- (16.034,-2.086) -- (16.400,-2.119)}
\def\profidhi{(8.350,-1.653) -- (8.716,-1.657) -- (9.082,-1.578) -- (9.448,-1.445) -- (9.814,-1.436) -- (10.180,-1.618) -- (10.545,-1.841) -- (10.911,-1.935) -- (11.277,-1.936) -- (11.643,-1.998) -- (12.009,-2.145) -- (12.375,-2.212) -- (12.741,-2.082) -- (13.107,-1.856) -- (13.473,-1.725) -- (13.839,-1.712) -- (14.205,-1.668) -- (14.570,-1.520) -- (14.936,-1.402) -- (15.302,-1.472) -- (15.668,-1.684) -- (16.034,-1.846) -- (16.400,-1.879)}
\def\profidlo{(8.350,-2.133) -- (8.716,-2.137) -- (9.082,-2.058) -- (9.448,-1.925) -- (9.814,-1.916) -- (10.180,-2.098) -- (10.545,-2.321) -- (10.911,-2.415) -- (11.277,-2.416) -- (11.643,-2.478) -- (12.009,-2.625) -- (12.375,-2.692) -- (12.741,-2.562) -- (13.107,-2.336) -- (13.473,-2.205) -- (13.839,-2.192) -- (14.205,-2.148) -- (14.570,-2.000) -- (14.936,-1.882) -- (15.302,-1.952) -- (15.668,-2.164) -- (16.034,-2.326) -- (16.400,-2.359)}
\def\profors{(8.350,-1.893) -- (8.716,-1.897) -- (9.082,-1.818) -- (9.448,-1.685) -- (9.814,-1.676) -- (10.180,-1.846) -- (10.545,-2.114) -- (10.911,-2.262) -- (11.277,-2.226) -- (11.643,-2.300) -- (12.009,-2.495) -- (12.375,-2.575) -- (12.741,-2.453) -- (13.107,-2.213) -- (13.473,-2.124) -- (13.839,-2.148) -- (14.205,-2.061) -- (14.570,-1.955) -- (14.936,-1.811) -- (15.302,-1.915) -- (15.668,-2.127) -- (16.034,-2.351) -- (16.400,-2.371) (8.350,-1.893) -- (8.716,-1.897) -- (9.082,-1.818) -- (9.448,-1.685) -- (9.814,-1.676) -- (10.180,-1.875) -- (10.545,-2.120) -- (10.911,-2.185) -- (11.277,-2.233) -- (11.643,-2.299) -- (12.009,-2.468) -- (12.375,-2.607) -- (12.741,-2.441) -- (13.107,-2.255) -- (13.473,-2.168) -- (13.839,-2.130) -- (14.205,-2.118) -- (14.570,-1.994) -- (14.936,-1.892) -- (15.302,-1.943) -- (15.668,-2.211) -- (16.034,-2.427) -- (16.400,-2.395) (8.350,-1.893) -- (8.716,-1.897) -- (9.082,-1.818) -- (9.448,-1.685) -- (9.814,-1.676) -- (10.180,-1.843) -- (10.545,-2.122) -- (10.911,-2.251) -- (11.277,-2.249) -- (11.643,-2.356) -- (12.009,-2.500) -- (12.375,-2.609) -- (12.741,-2.520) -- (13.107,-2.252) -- (13.473,-2.166) -- (13.839,-2.153) -- (14.205,-2.146) -- (14.570,-1.980) -- (14.936,-1.899) -- (15.302,-2.000) -- (15.668,-2.263) -- (16.034,-2.451) -- (16.400,-2.433) (8.350,-1.893) -- (8.716,-1.897) -- (9.082,-1.818) -- (9.448,-1.685) -- (9.814,-1.676) -- (10.180,-1.916) -- (10.545,-2.143) -- (10.911,-2.230) -- (11.277,-2.251) -- (11.643,-2.338) -- (12.009,-2.557) -- (12.375,-2.589) -- (12.741,-2.484) -- (13.107,-2.298) -- (13.473,-2.175) -- (13.839,-2.181) -- (14.205,-2.131) -- (14.570,-2.037) -- (14.936,-1.928) -- (15.302,-2.065) -- (15.668,-2.284) -- (16.034,-2.448) -- (16.400,-2.522) (8.350,-1.893) -- (8.716,-1.897) -- (9.082,-1.818) -- (9.448,-1.685) -- (9.814,-1.676) -- (10.180,-1.891) -- (10.545,-2.078) -- (10.911,-2.186) -- (11.277,-2.260) -- (11.643,-2.328) -- (12.009,-2.528) -- (12.375,-2.679) -- (12.741,-2.483) -- (13.107,-2.286) -- (13.473,-2.202) -- (13.839,-2.220) -- (14.205,-2.180) -- (14.570,-2.054) -- (14.936,-1.986) -- (15.302,-2.073) -- (15.668,-2.264) -- (16.034,-2.477) -- (16.400,-2.536) (8.350,-1.893) -- (8.716,-1.897) -- (9.082,-1.818) -- (9.448,-1.685) -- (9.814,-1.676) -- (10.180,-1.886) -- (10.545,-2.113) -- (10.911,-2.247) -- (11.277,-2.218) -- (11.643,-2.330) -- (12.009,-2.545) -- (12.375,-2.565) -- (12.741,-2.545) -- (13.107,-2.270) -- (13.473,-2.236) -- (13.839,-2.201) -- (14.205,-2.229) -- (14.570,-2.087) -- (14.936,-1.946) -- (15.302,-2.121) -- (15.668,-2.364) -- (16.034,-2.507) -- (16.400,-2.559) (8.350,-1.893) -- (8.716,-1.897) -- (9.082,-1.818) -- (9.448,-1.685) -- (9.814,-1.676) -- (10.180,-1.862) -- (10.545,-2.116) -- (10.911,-2.219) -- (11.277,-2.319) -- (11.643,-2.367) -- (12.009,-2.573) -- (12.375,-2.639) -- (12.741,-2.515) -- (13.107,-2.327) -- (13.473,-2.216) -- (13.839,-2.245) -- (14.205,-2.217) -- (14.570,-2.098) -- (14.936,-1.975) -- (15.302,-2.089) -- (15.668,-2.399) -- (16.034,-2.520) -- (16.400,-2.569)}
\def\profood{(8.350,-1.785) -- (8.716,-1.789) -- (9.082,-1.710) -- (9.448,-1.704) -- (9.814,-1.706) -- (10.180,-1.874) -- (10.545,-2.213) -- (10.911,-2.307) -- (11.277,-2.337) -- (11.643,-2.396) -- (12.009,-2.612) -- (12.375,-2.671) -- (12.741,-2.580) -- (13.107,-2.346) -- (13.473,-2.238) -- (13.839,-2.343) -- (14.205,-2.259) -- (14.570,-2.166) -- (14.936,-2.063) -- (15.302,-2.145) -- (15.668,-2.381) -- (16.034,-2.610) -- (16.400,-2.636)}
\def\profxA{16.400}\def\profyA{-2.119}\def\profyAor{-2.484}\def\profyAood{-2.636}
\def\profxone{8.350}

% ================================================================ (a)
\node[ttl] at (0,0.3) {\textbf{(a)}\ 向きの変え方は,そのままを除いて7つ};
\foreach \f/\x/\y/\t in {o0/0.7/-0.85/{}, vflip/2.35/-0.85/上下反転, hflip/4.0/-0.85/左右反転, r180/5.65/-0.85/180度回転,
                         r90/0.7/-2.75/90度回転, r270/2.35/-2.75/270度回転, tr/4.0/-2.75/{対角線で\\反転}, atr/5.65/-2.75/{もう1本の\\対角線で反転}}
  {\node[inner sep=0, draw=black!50, line width=0.4pt] at (\x,\y) {\includegraphics[width=1.15cm]{img/b_\f.png}};
   \node[sm, anchor=north] at (\x,{\y-0.64}) {\t};}
\draw[navyc, line width=1.3pt] (0.105,-0.255) rectangle (1.295,-1.445);
\node[sm, anchor=north, text=navyc] at (0.7,-1.49) {そのまま};
\node[sm, anchor=north west, align=left] at (0.0,-4.4) {そのまま以外の7つが $\mathbf{G}_{1}\vec{x},\dots,\mathbf{G}_{7}\vec{x}$(画素の並べ替え)};
% ================================================================ (b)
\draw[->, line width=0.9pt] (6.75,-1.8) -- (7.65,-1.8);
\node[sm, anchor=south] at (7.2,-1.75) {U-Netに};
\node[sm, anchor=north] at (7.2,-1.85) {通す};
\node[ttl] at (7.9,0.3) {\textbf{(b)}\ ブロックごとの振幅の平均(模式図)};
\draw[->, black!80] (8.05,-3.45) -- (16.95,-3.45);
\draw[->, black!80] (8.05,-3.45) -- (8.05,-0.3);
\node[lab, anchor=north] at (12.2,-3.5) {ブロック $h$};
\node[lab, anchor=west] at (8.15,-0.38) {振幅 $a_{h}$};
\node[sm, anchor=north] at (\profxone,-3.5) {1};
\node[sm, anchor=north, text=redc] at (\profxA,-3.5) {A};
\draw[black!45, densely dotted] (\profxA,-3.45) -- (\profxA,-0.6);
% ID band and mean
\draw[navyc!70, densely dotted, line width=0.5pt] \profidhi;
\draw[navyc!70, densely dotted, line width=0.5pt] \profidlo;
\draw[black!40, line width=0.45pt] \profors;
\draw[redc, line width=1.0pt] \profood;
\draw[navyc, line width=1.2pt] \profid;
\node[sm, text=navyc, anchor=south west, align=left] at (8.4,-1.2) {分布内の画像(そのまま). 点線：$\pm1$ 標準偏差};
\node[sm, text=black!60, anchor=north west] at (11.0,-3.0) {向きを変えた画像(7つの向き)};
\node[sm, text=redc, anchor=north west] at (8.65,-2.6) {分布外の画像};
\draw[->, line width=0.8pt] ({\profxA+0.25},\profyA) -- ({\profxA+0.25},\profyAood);
\node[sm, anchor=west] at ({\profxA+0.3},{(\profyA+\profyAood)/2}) {下がる};
\end{tikzpicture}
